\chapter{APRENDIZADO DE MÁQUINA}

Este capítulo apresenta os fundamentos do aprendizado de máquina supervisionado e os algoritmos centrais para os experimentos deste trabalho. São discutidos o paradigma supervisionado e os principais algoritmos utilizados em detecção de intrusão, a arquitetura e as propriedades do \textit{Random Forest} (com destaque para a técnica MDI de seleção de características), e as métricas de avaliação de classificadores em conjuntos de dados desbalanceados.

\section{Fundamentos do Aprendizado de Máquina Supervisionado}

O Aprendizado de Máquina (\textit{Machine Learning} - ML) é uma abordagem que confere aos sistemas a capacidade de aprender e melhorar a partir de dados, sem que sejam explicitamente programados para cada situação \cite{sbrc_minicursos}. No contexto da segurança de redes, essa capacidade é especialmente relevante: os algoritmos de ML conseguem identificar padrões complexos no tráfego que escapam às regras fixas dos sistemas tradicionais \cite{santos2023ercemapi}.

No aprendizado supervisionado, o modelo é treinado com um conjunto de dados cujas amostras já possuem rótulos de classe. Cada amostra é representada por um vetor de atributos (características) e associada a uma classe de saída, como tráfego benigno ou um tipo específico de ataque. A partir dessas associações, o modelo aprende a mapear características para classes e, após o treinamento, aplica esse mapeamento a dados inéditos \cite{sbrc_minicursos}. Em contraste, algoritmos não supervisionados inferem agrupamentos nos dados sem depender de rótulos, o que os torna adequados a cenários onde a anotação manual é custosa ou inviável \cite{santos2023ercemapi}. Abordagens semissupervisionadas combinam uma pequena quantidade de dados rotulados com um volume maior de dados não rotulados, aproveitando o melhor das duas estratégias.

Entre os algoritmos supervisionados mais utilizados em sistemas de detecção de intrusão estão as Árvores de Decisão, o \textit{Random Forest}, as Máquinas de Vetores de Suporte (SVM), o $k$-vizinhos mais próximos ($k$NN), as Redes Neurais Artificiais (ANN) e o \textit{XGBoost} \cite{santos2023ercemapi}. A escolha do algoritmo depende das características do problema, incluindo o tamanho do conjunto de dados, o número de classes, o grau de desequilíbrio entre elas e os requisitos de interpretabilidade do modelo.

\section{Árvores de Decisão e Random Forest}

Uma árvore de decisão é uma estrutura hierárquica de regras do tipo \textit{if-then-else} que divide recursivamente o espaço de atributos. Em cada nó interno, um atributo é selecionado com base em algum critério de impureza ou de ganho de informação, e o conjunto de amostras é dividido em subconjuntos para os nós filhos. As folhas da árvore correspondem às classes de saída. Três implementações clássicas diferem principalmente nos critérios de divisão: o ID3 maximiza o ganho de informação e minimiza a entropia; o CART realiza divisões binárias usando a impureza de Gini; e o C4.5, evolução do ID3, suporta atributos categóricos e numéricos e resolve o problema de sobreajuste com técnicas de poda \cite{sbrc_minicursos}. A principal vantagem das árvores de decisão é a interpretabilidade, pois as regras resultantes podem ser exibidas graficamente e inspecionadas por analistas.

O \textit{Random Forest} é um algoritmo baseado em conjunto (\textit{ensemble}) que cria múltiplas árvores de decisão durante o treinamento. Para cada árvore, um subconjunto aleatório das amostras é selecionado por reamostragem com reposição (\textit{bootstrap}) e um subconjunto aleatório dos atributos é considerado em cada divisão. A classificação final é determinada por votação majoritária entre todas as árvores \cite{santos2023ercemapi}. Essa combinação de aleatoriedade na seleção de amostras e atributos reduz a correlação entre as árvores individuais e torna o modelo mais robusto ao sobreajuste quando comparado a uma única árvore de decisão.

Uma propriedade relevante do \textit{Random Forest} para aplicações em detecção de intrusão é a capacidade de calcular a importância de cada atributo durante o treinamento. A métrica de Diminuição Média da Impureza (\textit{Mean Decrease in Impurity} - MDI) quantifica, para cada atributo, o quanto ele contribui para reduzir a impureza nos nós das árvores ao longo de toda a floresta. Atributos com MDI elevado são os mais discriminativos para separar as classes \cite{sbrc_minicursos}. Esse mecanismo permite usar o \textit{Random Forest} como um passo de seleção de atributos antes do treinamento de outros modelos, reduzindo a dimensionalidade do conjunto de dados sem a necessidade de um método externo de seleção.

\section{Avaliação de Modelos e Seleção de Características}

A avaliação do desempenho de um classificador parte da matriz de confusão, que registra, para cada classe, a quantidade de Verdadeiros Positivos (VP), Falsos Positivos (FP), Verdadeiros Negativos (VN) e Falsos Negativos (FN). A partir dessas contagens, derivam-se as principais métricas de avaliação \cite{sbrc_minicursos}. A acurácia representa a fração total de amostras classificadas corretamente, mas em conjuntos com classes muito desbalanceadas não é um indicador confiável, pois um modelo que classifica tudo como a classe majoritária pode atingir alta acurácia sem detectar nenhum ataque. A precisão mede a fração de alertas gerados que correspondem a ameaças reais, enquanto a revocação (\textit{recall} ou sensibilidade) mede a fração de ameaças reais que foram efetivamente detectadas. O F1-score é a média harmônica entre precisão e revocação e equilibra as duas métricas em um único valor. Em problemas multiclasse, o F1 macro calcula o F1 de cada classe individualmente e tira a média sem ponderação pelo tamanho das classes, conferindo peso igual a classes raras e frequentes.

O sobreajuste (\textit{overfitting}) ocorre quando o modelo aprende as particularidades e ruídos do conjunto de treinamento a ponto de prejudicar a capacidade de generalização para dados novos. O subajuste (\textit{underfitting}), por sua vez, resulta de modelos com baixa capacidade expressiva que não capturam os padrões relevantes \cite{sbrc_minicursos}.

A seleção de características busca identificar o subconjunto de atributos mais informativo para a tarefa de classificação. Métodos de filtragem avaliam os atributos de forma independente do classificador, utilizando heurísticas como o ganho de informação ou a correlação. Métodos de empacotamento avaliam subconjuntos de atributos usando o desempenho do classificador como critério. Métodos embarcados realizam a seleção como parte integrante do treinamento do modelo \cite{sbrc_minicursos}. O MDI do \textit{Random Forest} é um exemplo de método embarcado: a importância de cada atributo é calculada naturalmente durante a construção das árvores, sem custo computacional adicional. Em \textit{datasets} de tráfego de rede com dezenas ou centenas de atributos, a redução de dimensionalidade obtida por esse método diminui o custo de treinamento e inferência e pode atenuar o sobreajuste em modelos mais complexos \cite{medeiros2019sbrc}. Com os conceitos de classificação supervisionada e seleção de características estabelecidos, o capítulo seguinte aborda o segundo pilar teórico dos experimentos: as Redes Neurais Profundas e a arquitetura de \textit{Autoencoder} utilizada na abordagem não supervisionada.
